\chapter{Abstract}

Artificial-intelligence drug discovery (AIDD) now generates more candidate targets, molecules and protocols than any single researcher can triage, yet the tools that produce them are disjoint and the reasoning that connects them is rarely auditable. We present MOCHI (Multi-agent Orchestrator for Chemical and biomedical Intelligence), a self-evolving multi-agent system in which a manager, a developer, a critic and a tool-creation agent collaborate over a shared ocean of 46 scientific tools. MOCHI is deliberately scenario-agnostic: four interchangeable stage programs cover (i) downstream effector-target discovery, (ii) structure-guided enzyme variant design, (iii) multi-step retrosynthesis planning and (iv) compilation of wet-lab protocols into validated self-driving-laboratory (SDL) instructions. A deterministic critic audit, a provenance record of every database call and a dual offline/online engine make every run reproducible and inspectable.

On four real tasks MOCHI recovers a Wnt/Notch receptor-class shortlist that motivated the experimental validation of the Wnt co-receptor LRP6, ranks the experimentally validated best strictosidine-synthase variant first out of 2,773 candidates, plans a two-step synthesis of lidocaine and a three-step synthesis of benorilate (all eight benchmark drugs solved), and compiles six clean protocols while catching five of five deliberately corrupted ones. Ablations show that dynamic tool creation is required for the target task and that the experimental prior is required to place the true best variant first. The run artefacts and the command that rebuilds every table and figure are available at \url{https://github.com/simmzx/mochi}, and a single command reproduces the whole suite.

\Keywords{multi-agent systems, AI drug discovery, retrosynthesis, enzyme engineering, laboratory automation, reproducibility}
